\begin{abstract}
Questo studio mostra come i modelli di machine learning possano essere utilizzati per aiutare ad effettuare stime precise del deficit di pressione di vapore (VPD), valore fondamentale per la gestione delle risorse idriche in agricoltura soprattutto in zone del mondo che soffrono di stress idrico, il VPD è uno dei parametri dell'ET che ha un impatto significativo sul suo calcolo.\\
Questo studio stima il VPD utilizzando un modello basato sull'apprendimento d'insieme in tutte le province siciliane (Messina, Catania, Palermo, Caltanissetta, Agrigento, Siracusa,Ragusa,Trapani, Enna). Per eseguire le stime sono stati utilizzati otto algoritmi di machine learning: Linear Regression (LR), Additive regression trees (AR), Random subspace(RSS), Random Forest(RF), Reduced Error Pruning Tree (REPTree), l'algoritmo M5 di Quinlan (M5P), Long Short-Term Memory (LSTM) e Gated Recurrent Unit (GRU); I dati sono stati ottenuti dalla rianalisi sui dati medi mensili ERA5 relativi a singoli livelli, prendendo i dati climatici rilevati tra il 1958 ed il 2023.\\
Sono stati creati due diversi tipi di split del dataset tra train e test set, nella prima versione i dati climatici sono stati divisi prendendo i valori dell'arco temporale che va dal 1958 al 2011 per il train set e quelli che vanno dal 2011 al 2025 per il test set, per verificare le prestazioni dei modelli su regioni che il modello ha già conosciuto durante la fase di training, successivamente per testare le prestazioni dei modelli su province mai provate in fase di training è stato effettuato uno split spaziale creando 10 coppie diverse di province i cui valori sarebbero, di volta in volta, esclusi dalla fase di training per poi venire usate solo nella fase di testing.\\
Per valutare le prestazioni del modello sono state utilizzate cinque misure statistiche: il coefficiente di correlazione(CC), l'errore assoluto medio (MAE), l'errore quadratico medio (RMSE), l'errore assoluto relativo (RAE) e l'errore quadratico relativo (RRSE).

\end{abstract}